\chapter*{Core Course Outcomes \& Bloom's Taxonomy Guide}
\addcontentsline{toc}{chapter}{Core Course Outcomes \& Bloom's Taxonomy Guide}

\begin{itemize}[leftmargin=1.5em, itemsep=4pt]
    \item \textbf{CO1}: Understand procedural Python syntax, variable binding, numeric primitives, and operator precedence rules.
    \item \textbf{CO2}: Construct robust control-flow algorithms utilizing conditional branching, structured loops, and recursive function designs.
    \item \textbf{CO3}: Select, manipulate, and optimize composite linear and non-linear data structures (strings, lists, tuples, and dictionaries).
    \item \textbf{CO4}: Leverage NumPy N-dimensional arrays, vectorization, and multi-dimensional broadcasting for high-performance numerical computing.
    \item \textbf{CO5}: Perform structured tabular data processing and aggregation using Pandas, and produce publication-grade statistical charts with Matplotlib.
    \item \textbf{CO6}: Implement integrated Python pipelines connecting digital image processing (OpenCV), interactive GUI loops (Pygame), and automated web scraping.
\end{itemize}

\vspace{5mm}
\noindent
\textbf{Bloom's Taxonomy Levels Utilized in Question Marking:}
\begin{enumerate}[label=\textbf{L\arabic*:}, leftmargin=1.5em, itemsep=2pt]
    \item \textbf{Remember:} Stating definitions, keywords, method signatures, and default behaviors.
    \item \textbf{Understand:} Explaining control flow, tracing variable scope (LEGB), and predicting output of code snippets.
    \item \textbf{Apply:} Implementing algorithmic logic, correcting syntax traps, and computing transformations.
    \item \textbf{Analyze:} Formulating time and space complexity $\mathcal{O}(n)$, diagnosing memory mutations, and optimizing vectorized operations.
\end{enumerate}

\newpage
